\hypertarget{const-volatile}{%
\chapter{CONST \& VOLATILE}\label{const-volatile}}

\hypertarget{const-correctness}{%
\section{11.1 Const Correctness}\label{const-correctness}}

\begin{lstlisting}[language={C++}]
#include <iostream>
using namespace std;

int main() {
    int x = 10;
    
    // const variable
    const int constant = 5;
    // constant = 10;  // ERROR
    
    // pointer to const
    const int* ptr1 = &x;
    // *ptr1 = 20;  // ERROR
    ptr1 = &constant;  // OK
    
    // const pointer
    int* const ptr2 = &x;
    *ptr2 = 20;  // OK
    // ptr2 = &constant;  // ERROR
    
    // const reference
    const int& ref = x;
    // ref = 20;  // ERROR
    
    cout << x << endl;
    cout << *ptr1 << endl;
    cout << *ptr2 << endl;
    cout << ref << endl;
    
    return 0;
}
\end{lstlisting}

\hypertarget{volatile-keyword}{%
\section{11.2 Volatile Keyword}\label{volatile-keyword}}

\begin{lstlisting}[language={C++}]
#include <iostream>
using namespace std;

int main() {
    // volatile - tells compiler value may change unexpectedly
    volatile int sensor_reading = 0;  // From hardware
    
    // Compiler won't optimize away reads
    while (sensor_reading < 100) {
        // Check actual value each time, not cached
    }
    
    // Common use: hardware registers
    volatile int* hardware_register = (volatile int*)0x1000;
    
    // Each access reads from actual location
    int val1 = *hardware_register;
    int val2 = *hardware_register;
    
    // Without volatile, compiler might optimize one read
    
    return 0;
}
\end{lstlisting}

\begin{center}\rule{0.5\linewidth}{0.5pt}\end{center}
